% ======================================================================
% Chapter 22: Quantum Field Theory at the Frontier
% (book pages 781-823; PDF pages 804-823) -- The Epilogue
% ======================================================================
\chapter{前沿的量子场论}

在本教科书中，我们考察了量子场论中最重要的思想。从融合相对论、量子力学与场的那些基本概念出发，我们构建了一个精巧的结构，它包含了诸如耦合常数重整化和非阿贝尔规范场这样非凡的要素。我们在许多地方都看到，这一结构中那些奇异而抽象的要素实际上与观测相交汇，甚至为基本粒子行为的意外方面提供了解释。

在我们研究的过程中，我们到达了一个对基本粒子的强、弱和电磁相互作用的完整理论。这一理论的每一个要素都被描述为一种量子场论，而这些量子场论被证明具有非常相似的结构，都是与费米子耦合的规范理论。在我们讨论的各个地方，我们都注意到这些理论已通过了严格的定量实验检验。我们没有篇幅来描述那些增进我们对这些理论信心的种类繁多的实验，但如今几乎所有粒子物理学家都认为这个 $SU(3)\times SU(2)\times U(1)$ 规范理论已经确立。事实上，这些人中的大多数都略带轻视地称这个理论为「标准模型」。

尽管支持标准模型的最佳数据来自过去五年的实验，但其背后的思想要古老得多。本书描述的大部分理论进展都在 20 世纪 70 年代完成，距离当今物理学前沿已有一代人之遥。但这并不意味着量子场论与那前沿无关，正如量子力学和电动力学在多年的探索之后并未失去其意义一样。恰恰相反，基本粒子理论——如同物理学中其他依赖连续介质中量子涨落的领域一样——仍然蕴含着深刻的奥秘，而量子场论仍是探索这些问题的主要工具。

在这最后一章中，我们将闪回到当下，讨论量子场论与基本相互作用物理学中当前问题的相关性。我们将提出在我们看来基本粒子物理学中尚未解决的重要问题，并描述量子场论如何被用来直面这些问题。这些应用中有许多涉及量子场论中超出本书范围的方面。这包括在微扰论无法企及的领域中运用量子场论，以及运用量子场论来探索具有更高自旋和局域对称性的理论的特殊性质。在这些领域中，我们的讨论将主要是定性的，但我们会给出一些参考文献，为这些主题各自提供切入点。

显而易见，我们在最后一章中的讨论将表达我们个人的观点，而绝不代表量子场论专家们的共识。此外，在一个快速发展的研究领域中，任何关于「当前问题」的汇集都会很快过时。事实上，我们希望本书的读者能够通过解决我们在此着重指出的问题，很快使这一章变得过时。

\section{强强相互作用}

我们对强相互作用的讨论有一个矛盾之处：我们所有具体的结果都是通过假定这些相互作用很弱而得到的。我们论证过，在大动量转移下，由于渐近自由，这个假设实际上成立。尽管如此，我们把关于强相互作用粒子最显而易见的问题——例如它们的质量以及低能相互作用——留在了被排除在我们分析之外的神秘区域，这仍令人不安。

要在强相互作用很强的区域中运用 QCD，我们需要回答三个问题：第一，我们如何描述把夸克束缚成强子的力？第二，对这些力所束缚的夸克-反夸克系统和三夸克系统，什么才是恰当的描述？最后，我们如何利用这些束缚态来计算散射振幅和流矩阵元？

目前，还没有从 QCD 拉格朗日量出发导出夸克之间完整力的途径。显式的计算只能在弱耦合和强耦合的极限下进行。在弱耦合极限下，结果是具有渐近自由耦合常数的库仑势。另一方面，强耦合极限给出一个线性势，它以我们在第~17.1 节末尾描述过但未推导的方式把颜色禁闭起来。这一结果的推导相当不寻常，并引入了一套新的数学方法。

到目前为止，本书尚未讨论量子场论的强耦合近似。原因很简单：在耦合 $g^2$ 很大的量子场论中，基本粒子或其束缚态通常获得随 $g^2$ 增长的质量。当 $g^2\to\infty$ 时，这些质量变得与截断 $\Lambda$ 相当，该场论便不再具有局域连续统描述。

Wilson 提出以激进的方式解决这个问题，即用离散间隔点的格点取代时空。这样的格点在欧几里得时空中最容易想象，因此我们可以用格点上场的泛函积分来逼近欧几里得格林函数。这样的理论可以具有良好定义的强耦合极限。这类理论与磁系统的格点模型非常相似。

事实上，我们可以利用第~13 章的概念来理解这种量子场论的构造。在每个格点处具有涨落自旋变量的格点理论，在宏观上由具有底层自旋变量对称性的标量场量子场论来描述。通常，量子场论的强耦合极限对应磁体的高温极限，此时关联长度远小于格点间距。减小耦合常数对应于降低温度。最终，耦合常数趋近重整化群的不动点，人们可以利用这一不动点来定义格点泛函积分中格点间距趋于零的极限。

要建立强相互作用的格点模型，人们需要在离散格点上找到一组在宏观上对应非阿贝尔规范场的变量。Wilson 提出，这样一个理论的基本变量应当是从一个格点顶点 $v_1$ 到相邻顶点 $v_2$ 的线元，
\begin{equation}
  U(v_1, v_2) ,
  \label{eq:22.1}
\end{equation}
其中每个 $U$ 都是规范群 $G$ 的一个元素。要构造规范群为 $G$ 的格点规范理论，人们应当对格点中每一条链环对一个有限群变换 $U$ 作积分。把这些 $U$ 矩阵沿闭合路径取乘积，人们就可以构造规范不变的观测量，正如我们在第~15.3 节所做的那样。合适的拉格朗日量也可以构造为格点初等闭合环上 $U$ 矩阵的规范不变乘积之和。%
\footnote{这一构造由 K.~Wilson 提出，\emph{Phys.\ Rev.\ D} \textbf{10}, 2445 (1974)。M.~Creutz 的 \emph{Quarks, Gluons, and Lattices}（Cambridge University Press, Cambridge, 1983）对这一构造作了教学式的介绍。}

在自旋系统中，高温相的定义性性质是关联的指数衰减
\begin{equation}
  \langle s(0)\cdot s(x)\rangle \sim \exp[-|x|/\xi]
  \label{eq:22.2}
\end{equation}
当 $|x|\to\infty$ 时。闭合路径 $P$ 周围 $U$ 矩阵的规范不变关联函数的类似性质是
\begin{equation}
  \bigl\langle \mathrm{tr}\prod_P U \bigr\rangle \sim e^{-\sigma A},
  \label{eq:22.3}
\end{equation}
其中 $A$ 是路径所张开的面积。这种行为实际上可以在 Wilson 格点规范理论的强耦合展开中被显式看到。相距为 $R$、持续欧几里得时间 $T$ 的一对色源，由宽度为 $R$、长度为 $T$ 的大矩形圈表示。对于这样的圈，我们可以把 \eqref{eq:22.3} 的结果与欧几里得时间中激发态能量的表达式比较，
\begin{equation}
  \langle \exp[-H_E T] \rangle \sim \exp[-E\,T] .
  \label{eq:22.4}
\end{equation}
于是我们看到，在强耦合极限下，静规范荷源之间以势能
\begin{equation}
  V(R) \sim \sigma R
  \label{eq:22.5}
\end{equation}
相互吸引，在足够大的 $R$ 处成立，其中 $\sigma$ 是弦张力。类似地，当人们把动力学夸克引入格点规范理论并在强耦合极限下研究其性质时，色源大距离分离的组态在欧几里得泛函积分中会被 \eqref{eq:22.3} 形式的因子所压低。于是强耦合极限预言夸克被永久禁闭到色单态束缚态之中。

我们刚才给出的论证同样适用于基于阿贝尔或非阿贝尔对称群的规范理论。但非阿贝尔规范理论还拥有渐近自由这一重要的额外性质。在此语境下，这意味着短距离处弱耦合的理论会流向长距离处强耦合的理论。如果我们设想像第~12.1 节描述的那样积掉短距离自由度，并且如果 $\beta$ 函数没有零点或其他阻碍重整化群流动的障碍，我们最终应当到达一个强耦合展开是良好近似的有效理论。因此，在非阿贝尔规范理论的特殊情形下，渐近自由使我们能够把基于自由夸克与胶子的短距离图像，与基于色禁闭的长距离图像联系起来。

如果我们所描述的强耦合图像能在连续时空中导出描述被永久禁闭的夸克与反夸克运动的数学方程，那将是极好的。许多作者曾试图写出这样的方程，他们设想 Wilson 圈关联函数 \eqref{eq:22.3} 的面积压低源自一张张开该圈的物理曲面。对于与色源相关联的大矩形圈，这张曲面在物理上可以解释为从一个源延伸到另一个源的色电通量线（如图~\ref{fig:17.1} 所示），并在欧几里得时间内扫过。在欧几里得时间的某一时刻，这张曲面可以理想化为一种抽象的、通常称为\emph{弦}的一维激发。遗憾的是，理想化弦的量子性质被证明非常复杂，因为弦的每一个小单元都必须被视为独立的量子自由度。唯一实际被解出的弦方程组都具有奇异的特征，包括不受欢迎的无质量粒子。到目前为止，还没有人成功写出一个对夸克束缚态定量计算有用的、描述夸克禁闭弦的方程。%
\footnote{对于从涉及 Wilson 圈与弦的图像处理色禁闭的一种途径，见 A.~A.~Migdal, \emph{Phys.\ Repts.} \textbf{102}, 199 (1983)。}

然而，非阿贝尔规范理论的格点正则化提示了强相互作用理论中定量计算的另一种途径。通过用具有非零格点间距和有限时空体积的格点规范理论来逼近 QCD，我们把泛函积分约化为有限个有界积分，即对格点中有限数目的每一条链环对 $SU(3)$ 群矩阵作积分。例如，大小为 $20^4$ 的格点允许格点间距小于强子的大小，同时格点的整体大小远大于强子半径。于是，人们可以通过蒙特卡罗方法数值计算这些积分，从而计算关联函数。由于具有有限格点间距的泛函积分与原始具有零格点间距的泛函积分通过积掉短距离自由度而相关联，格点近似可以通过微扰地计算短距离效应来系统改进，并利用渐近自由来为弱耦合分析辩护。%
\footnote{数值格点规范理论的入门介绍见 \emph{From Actions to Answers}, T.~DeGrand 与 D.~Toussaint 编（World Scientific, Singapore, 1990）。}

这种数值方法现已成为强子物理定量计算的主要理论工具。目前，这一方法给出低激发介子和重子质量的精度为 10--20\%；它还允许以 25\% 的水平计算强子的弱相互作用矩阵元。随着计算机变得更强大，这种数值方法可以达到更高的精度。

最终，询问这些非微扰数值计算是否与我们关于 QCD 微扰区域的精确知识相一致，将是饶有趣味的。在撰写本书之时，第一例这样的比较已经做出：我们在表~\ref{tab:17.1} 中列出了从 $J/\psi$ 和 $\Upsilon$ 谱学得到的 $\alpha_s$ 值。在这一计算中，实验确定的 $c\bar{c}$ 和 $b\bar{b}$ 束缚态质量与用格点正则化数值计算得到的值相比较。这些值的比较给出所要求的格点理论裸耦合常数，它可以按表中的约定转换为 $\alpha_s(m_Z)$ 的值。由此得到的 $\alpha_s(m_Z)$ 估计值确实与纯微扰的测定相当一致。

强子物理中非微扰计算的未来如何？一方面，我们期望看到数值格点方法的进一步发展。这些方法几乎尚未开始处理强子-强子散射和多粒子矩阵元的问题，而这似乎是未来一个重要方向。此外，这些方法最终应当在百分之几或更好的水平上提供对强子的工程式理解。另一方面，我们也希望看到一种对强子结构的定量连续统途径，其中动力学夸克与某种恰当的弦自由度相互作用。

\section{大统一及其悖论}

如果我们把关于 QCD 低能非微扰行为的问题搁置一旁，$SU(3)\times SU(2)\times U(1)$ 规范理论在我们实验上探测过的那些能量处，对基本粒子相互作用给出了表面完整的描述。但超出我们当前所及的范围会发生什么？这个理论需要修改吗，还是它在高得多的能量下仍然有效？

$SU(3)\times SU(2)\times U(1)$ 规范理论包含三个独立的规范耦合常数，而这些耦合的观测值对规范群较大的分量而言更大。这种格局可以用一个关于规范耦合在甚高能量下行为的大胆假设来解释。如果在某个非常大的能量标度处这三个耦合相等，那么 $SU(3)$ 和 $SU(2)$ 耦合的值由于其渐近自由的重整化群方程会在较小动量标度处增大，而 $U(1)$ 耦合的值会减小，从而在低能处产生观测到的耦合格局。

一个更大胆的假设是，三种规范对称性是单一大型对称群的子群，该群在甚高能标处自发破缺。这个更大对称群的最简单选择是 $SU(5)$。在那个理论中，$SU(3)\times SU(2)\times U(1)$ 的耦合常数与 $SU(5)$ 破缺标度处底层 $SU(5)$ 耦合有如下关系：
\begin{equation}
  g_3 = g_2 = \sqrt{\frac{5}{3}}\,g_1 .
  \label{eq:22.6}
\end{equation}
$SU(3)\times SU(2)\times U(1)$ 规范群被嵌入一个更大的单群的这一思想，被称为\emph{大统一}；以 $SU(5)$ 作为统一群的特定选择归功于 Georgi 与 Glashow。%
\footnote{大统一理论综述见 P.~Langacker, \emph{Phys.\ Repts.} \textbf{72}, 185 (1981)，以及 P.~Langacker, \emph{Phys.\ Repts.} \textbf{109}, 145 (1984)。} 可以证明，观测到的夸克和轻子恰好嵌入了 $SU(5)$ 的一个无反常手征表示；这一嵌入自然地解释了夸克的分数电荷。

在这个框架内，我们可以把三个耦合常数的值从 $m_Z$ 能标向外推到更高能量。这一外推的结果如图~22.1 中的实线所示。耦合常数确实在甚高能量处彼此靠近，尽管它们并未真正相交。图中的虚线显示的是用一组修正过的重整化群方程得到的演化，这将在第~22.4 节中解释；采用这一选择，三个耦合在当前的不确定性范围内精确相交。无论如何，耦合常数的演化发生在能量的对数标度上，因此耦合相交的能量只能以对数方式确定。利用标准模型对三个耦合外推的结果是量级为
\begin{equation}
  M_U \sim 10^{15}\ \mathrm{GeV} .
\end{equation}
的大统一标度。

大统一作出了若干戏剧性的预言。由于每一代夸克和轻子被放入 $SU(5)$ 的单一表示中，该理论预言夸克与轻子之间的跃迁。这些跃迁由新的规范玻色子——$X$ 和 $Y$ 玻色子——所介导，它们极其重，量级为 $M_U$。这些相互作用违反重子数，导致质子衰变。预言的质子寿命为
\begin{equation}
  \tau_p \sim \frac{M_U^4}{\alpha_U^2 m_p^5} \sim 10^{30}\ \mathrm{years} ,
  \label{eq:22.7}
\end{equation}
这比宇宙年龄大许多个数量级。没有观测到质子衰变，这对大统一标度施加了严格限制，并排除了最简单的 $SU(5)$ 模型。尽管如此，大统一仍然是一个重要而有吸引力的思想。

\subsection{规范层级问题}

大统一导致一个关于希格斯玻色子质量的严重概念问题。为了产生大小恰当的真空期望值，以给出观测到的 $W$ 和 $Z$ 玻色子质量，希格斯标量场必须获得一个负质量项，其大小为
\begin{equation}
  -\mu^2 \sim -(100\ \mathrm{GeV})^2 .
  \label{eq:22.8}
\end{equation}
不幸的是，标量场的 $(\mathrm{mass})^2$ 会接受相加性重整化。在截断标度为 $\Lambda$ 的理论中，只有当标量场的裸质量量级为 $-\Lambda^2$、并且这个值被辐射修正对消到 $-\mu^2$ 时，$\mu^2$ 才能远小于 $\Lambda^2$。如果我们设想我们关于自然的理论包含大统一的甚大标度，我们就必须认真对待这样的观念：在此讨论中 $\Lambda$ 应取的恰当值是 $10^{16}$ GeV 或更大。这似乎要求 $\mu^2$ 的重整化值发生戏剧性甚至离奇的对消。

我们在相变理论中遇到过这类情形。在零温下，铁磁体的自旋期望值通常与底层原子参数同量级。随着温度升高，或系统中某个其他变量被改变，磁化强度减小。最终，通过对温度的精细调节，我们可以到达系统趋近临界点的情形。在非常靠近该点的区域内，自旋场的期望值远小于由原子参数预言的值，系统由一个近似无质量的连续统标量场理论描述。

在统计力学中，这种轻标量场的图像之所以成立，是因为有一位实验者灵敏地调节旋钮。在弱相互作用理论中，显然没有人做这种精细调节，使希格斯玻色子的 $(\mathrm{mass})^2$ 的值比其自然值低 28 个或更多数量级。因此，为什么希格斯玻色子质量与大统一标度相比如此之小，是一个谜。粒子物理学家把这个问题称为\emph{规范层级问题}。

人们如何自然地安排一个希格斯场质量项，使其远小于基本相互作用的底层质量标度？一种可能的策略是安排基本拉格朗日量的一种对称性，它禁止希格斯玻色子质量项，并且只是非常弱地破缺。这一思想被证明极难实现。要构建这类理论，人们需要创建一个标量场理论，其中对希格斯玻色子质量的相加性辐射修正必须在微扰论的任何可预见的阶都相互对消。但希格斯质量项的形式非常简单，
\begin{equation}
  \Delta\mathcal{L} = \mu^2|\phi|^2 ,
  \label{eq:22.9}
\end{equation}
很难想象有任何原理能阻止这个项被辐射修正产生。对于具有这一性质的对称性有一个提议，但它需要引入一个被称为\emph{超对称}的深刻原理，它要求对基本物理学作深刻的修改。特别是，它要求大量新的基本粒子，其中一些的质量应低于 1000 GeV，处于下一代加速器所能及的范围之内。我们将在第~22.4 节进一步讨论这一可能性。

在这一讨论中，希格斯质量问题源于希格斯玻色子是基本粒子的假设。另一个已经在第~20.2 节末尾提出的观点是，希格斯玻色子是由一套新的相互作用束缚的复合态。这一思想也导致可观测的实验后果，因为这些新相互作用的质量标度必须接近弱相互作用质量标度。在这类理论的最简单形式中，希格斯扇区的对称性破缺以强相互作用的动力学手征对称性破缺为模型，后者我们在第~19.3 节讨论过。该理论所要求的新强相互作用导致一套质量约为 1000 GeV 的新粒子谱。%
\footnote{这些希格斯扇区模型（专家们称之为 technicolor 模型）的性质见 R.~Kaul, \emph{Rev.\ Mod.\ Phys.} \textbf{55}, 449 (1983)，以及 K.~D.~Lane, 载于 \emph{The Building Blocks of Creation}, S.~Raby 编（World Scientific, 1993）。} 因此，关于破缺 $SU(2)\times U(1)$ 的扇区性质的两种相互冲突的假设，都导致在未来加速器、甚至可能现在的加速器上可观测的新现象。

正如希格斯扇区的这两种不同理论对弱相互作用对称性 $SU(2)\times U(1)$ 为何应当自发破缺的问题给出完全不同的回答一样，它们对夸克和轻子质量起源的问题也给出完全不同的回答。在希格斯场是基本的模型中，夸克和轻子质量来自费米子与希格斯场的可重整耦合。这些耦合大概会是大统一的一部分，并且只能由明确涉及大统一标度的理论来预言。原则上，对这些耦合的了解可以给我们提供大统一细节的线索。即使希格斯场是复合的，我们也不能回避这样一个事实：夸克和轻子质量的产生要求 $SU(2)\times U(1)$ 的破缺。因此，这些质量项必定来自夸克和轻子与希格斯相互作用扇区的耦合。在这类模型中，导致夸克和轻子质量的相互作用必须在接近希格斯扇区强相互作用标度的能量处出现，并最终可能在实验上可观测。

从任一观点来看，为什么夸克和轻子谱跨越 5 个数量级——从 0.5 MeV 的电子到 175 GeV 的顶夸克——仍然是个谜。是什么产生编码在 CKM 矩阵中的夸克混合格局以及 $CP$ 破坏的幅度，也尚未被理解。即使有了对标准模型的详细证实，这些问题在今天看来也远未得到解决。

\subsection{宇宙学常数问题}

大统一的巨大质量标度还可能进入另一个物理量，它所构成的悖论甚至比希格斯玻色子质量更大。当我们在第~2.3 节首次量子化一个场时，我们发现自由标量场论中真空的能量密度从各种振荡模式的零点能获得一个无限的正贡献。在截断标度为 $\Lambda$ 时，这个零点能大致为
\begin{equation}
  \langle 0|H|0\rangle \sim \Lambda^4 .
  \label{eq:22.10}
\end{equation}
在我们讨论的许多其他地方，我们都发现了对真空能类似的大贡献。在自由费米子理论的量子化中，Dirac 海的填充导致真空能的一个具有类似紫外发散的向下偏移。自发对称性破缺给出真空能密度的有限但可能仍然很大的偏移，
\begin{equation}
  \Delta\langle 0|H|0\rangle \sim -cv^4 ,
  \label{eq:22.11}
\end{equation}
其中 $c$ 无量纲，$v$ 为场的真空期望值。弱相互作用 $SU(2)\times U(1)$ 对称性与强相互作用手征对称性的自发破缺，预期都会以这种方式使真空能密度发生偏移。

在基本粒子物理实验中，真空能的这种偏移是不可观测的。例如，实验测量的粒子质量是真空与 $H$ 的某些激发态之间的能量差，绝对真空能在这些差值的计算中会抵消。然而，有一种方式可能观测到绝对真空能，即通过真空能与引力的耦合。按照 Einstein，物质的能量-动量张量 $\Theta^{\mu\nu}$ 是引力场的源。真空能密度 $\langle 0|H|0\rangle = \Lambda$ 对这一源贡献一项
\begin{equation}
  \Theta^{\mu\nu} \supset \Lambda\,g^{\mu\nu} ,
  \label{eq:22.12}
\end{equation}
其中右边的第一项被减去以具有零真空期望值。真空能项具有 Einstein 的\emph{宇宙学常数}的形式，因此可能影响宇宙的膨胀。

事实上，对宇宙学膨胀的测量排除了大的宇宙学常数。当前的极限是
\begin{equation}
  \Lambda < 10^{-29}\ \mathrm{g/cm}^3 \sim (10^{-11}\ \mathrm{GeV})^4 .
  \label{eq:22.13}
\end{equation}
我们不理解为什么 $\Lambda$ 比粒子物理已知相变中产生的真空能偏移小这么多，而又比底层的场零点能小得更多。实验界 \eqref{eq:22.13} 与朴素直觉之间 $\Lambda$ 的出入是 120 个数量级！这个问题的解决方案可能来自许多来源之一。可能引力量子场论的形式体系要求把真空能从出现在 Einstein 引力方程中的能量-动量张量中减去。可能来自粒子物理或引力本身的新物理机制把总真空能设为零。也可能能量-动量的整体标度确实是有歧义的，由某个宇宙学边界条件来设定。此时此刻，所有这些可能性都只是猜测。我们唯一确定知道的是，量子场论与引力的统一不可能是一帆风顺的，我们当前的理解中还缺少某个重要的概念。%
\footnote{宇宙学常数问题以及各种不成功的解决方案的综述见 S.~Weinberg, \emph{Rev.\ Mod.\ Phys.} \textbf{61}, 1 (1989)。}

\section{量子场论中的精确解}

从大统一这一充满巨大前景与奥秘的思想，我们转向对模型量子场论的研究，这些模型简单到可以被精确求解。在整本书中，我们都强调了量子场论的内在复杂性，以及用微扰论替代精确知识的重要性。但存在各种各样的已知精确解的量子场论。在本节中，我们将描述其中一些，并回顾我们从它们获得的洞见。

在寻找量子场论模型的精确解时，没有理由把我们注意力限制在四维时空。事实上，我们已经见过具有类似重整化和短距离行为复杂性的二维理论的例子。同时，这些理论占据一维空间，它们的自由度可以被想象成一条链中的链环。这允许一些强有力的简化。

在我们第~19.1 节对二维轴反常的讨论中，我们证明了二维无质量 QED 的光子变为有质量玻色子。对这一理论的更细致考察表明，这个玻色子是一个无相互作用的粒子。该理论最初以费米子表述，通过库仑力相互作用。然而，可以把该理论精确改写为产生和湮灭费米子-反费米子对的标量场理论。启发式地说，在一维空间中沿光锥运动的一个粒子和一个反粒子不会分离，因此构成一个玻色子自由度。在一大类模型中，如此改写的玻色子理论是自由场理论。这类理论中一个显著的模型是\emph{Thirring 模型}，
\begin{equation}
  \mathcal{L} = \bar{\psi}i\slashed{\partial}\psi
    - \frac{1}{2}g(\bar{\psi}\gamma^\mu\psi)^2 ,
  \label{eq:22.14}
\end{equation}
在二维中。在这一模型中，用玻色子场替换费米子场得到一个自由场理论。利用这一场论，人们可以显式地计算费米子双线性算符的关联函数，并直接证明这些算符具有反常维度。用重整化群的语言说，该模型包含一条由耦合常数 $g$ 参数化的不动点直线。%
\footnote{这些模型的入门介绍见 S.~Coleman, \emph{Phys.\ Rev.\ D} \textbf{11}, 2088 (1975), \emph{Ann.\ Phys.} \textbf{101}, 239 (1976)。}

更一般的一类二维模型可以通过在哈密顿图像中把它们想象成耦合场算符的一维链来求解。这一方法的原型是磁体统计力学中的一个问题，即耦合自旋的一维链。考虑一条由 $N$ 个离散格点组成的长链，每个格点处有一个自旋-$1/2$ 系统。Pauli sigma 矩阵 $\bm{\sigma}_i$ 作用在格点 $i$ 处的二维 Hilbert 空间上。于是自旋链的哈密顿量为
\begin{equation}
  H = \sum_i \Bigl\{ -J\bm{\sigma}_i\cdot\bm{\sigma}_{i+1} \Bigr\} .
  \label{eq:22.15}
\end{equation}
由于
\begin{equation}
  \bm{\sigma}_i\cdot\bm{\sigma}_{i+1} = 2\bigl( \sigma^+_i\sigma^-_{i+1}
    + \sigma^-_i\sigma^+_{i+1} \bigr) + \sigma^z_i\sigma^z_{i+1} ,
  \label{eq:22.16}
\end{equation}
这个哈密顿量守恒上自旋的数目。所有自旋向下的态是哈密顿量的本征态，而一个自旋向上且具有确定动量的态也是本征态。1934 年，Bethe 分析了两个自旋向上的问题，并计算了它们的 $S$ 矩阵。随后他发现了令人惊奇的事实：通过把两自旋问题的 $S$ 矩阵相乘，他可以找到任意数目上自旋情况下哈密顿量的精确本征态。通过考虑 $N/2$ 个自旋向上，他找到了系统的基态。这一现在被称为\emph{Bethe 假设}的技巧，已被用来解决凝聚态物理和量子场论中各种各样的一维问题。例如，Andrei 和 Lowenstein 曾用这一技巧求解习题~11.3 中提出的 Gross-Neveu 模型，并证明该模型的谱具有渐近自由所预期的性质。%
\footnote{Bethe 假设及其推广的入门介绍见 N.~Andrei, K.~Furuya, 和 J.~H.~Lowenstein, \emph{Rev.\ Mod.\ Phys.} \textbf{55}, 331 (1983), L.~D.~Faddeev, 载于 \emph{Recent Advances in Field Theory and Statistical Mechanics}, J.~B.~Zuber 与 R.~Stora 编（North-Holland, Amsterdam, 1984），以及 R.~J.~Baxter, \emph{Exactly Solved Models in Statistical Mechanics}（Academic Press, London, 1982）。}

即使在无法对参数的所有值求解一个模型的情况下，有时也可以在二维模型包含无质量场的那些点处找到精确信息。众所周知，许多经典的二维偏微分方程可以利用复变函数技巧求解。例如，二维 Laplace 方程 $\nabla^2\phi = 0$ 对共形映射 $z\to w(z)$ 不变，其中 $z = x + iy$。具有无质量粒子的二维量子场论通常在经典层次上具有这种\emph{共形对称性}，尽管一般地它是反常的。然而，在特殊的系统中，这些反常消失，量子理论对共形映射不变。这些理论通常包含具有反常维度的算符，表明每一个这样的理论都是重整化群的一个新的、非平凡的不动点。理论的共形对称性可以用来计算这些反常维度。

作为这类理论的一个例子，考虑二维非线性 sigma 模型，其中基本场不是像我们在第~13.3 节讨论的那样是单位矢量，而是 Lie 群 $G$ 的一个幺正矩阵。该理论的拉格朗日量为
\begin{equation}
  \mathcal{L} = \frac{1}{4\lambda^2}\ \mathrm{tr}\bigl[ \partial_\mu U^{-1}
    \partial^\mu U \bigr] ,
  \label{eq:22.17}
\end{equation}
其中 $U$ 是在群 $G$ 中取值的场。与第~13.3 节的理论一样，该模型是渐近自由的。然而，Witten 已经证明，通过给这个拉格朗日量加上一个形式相当复杂、最早由 Wess 和 Zumino 写出的特定微扰项，人们可以找到一个具有明显 $G\times G$ 整体对称性的重整化群不动点。这个理论是共形不变的，所有算符关联函数都可以利用共形对称性来计算。%
\footnote{共形场论的入门介绍见 A.~A.~Belavin, A.~M.~Polyakov, 和 A.~B.~Zamolodchikov, \emph{Nucl.\ Phys.\ B} \textbf{241}, 333 (1984)，以及 J.~L.~Cardy, 载于 \emph{Phase Transitions and Critical Phenomena}, C.~Domb 与 J.~L.~Lebowitz 编（Academic Press, 1987）。}

对量子场论非微扰探索的一个结果是，人们意识到场论可以包含与原始场的量子没有简单关系的粒子态。在量子场论的弱耦合极限中，这样的新态可以作为经典场方程的新解出现。例如，考虑破缺对称性相中二维的 $\phi^4$ 理论。运动方程为
\begin{equation}
  \frac{\partial^2\phi}{\partial t^2} - \nabla^2\phi + \cdots = 0 ,
  \label{eq:22.18}
\end{equation}
它具有在两个简并真空之间插值的静态解，即\emph{扭结}。这些孤子态是经典解的量子力学对应物。类似的孤子激发，例如磁单极子和瞬子，在更现实的理论中起着重要作用。

\section{量子场论与引力}

量子场论的最终前沿是量子化引力本身的问题。引力相互作用在经典上由 Einstein 的广义相对论描述。在量子层次上，引力子——引力场的量子——是一个无质量的自旋-2 粒子。无质量自旋-2 场的理论受到规范不变性的高度约束，其方式与我们第~15 章研究的非阿贝尔规范理论所受的约束类似。

然而，引力的量子化遇到一个根本困难：该理论按幂次计数是不可重整的。引力耦合常数，即 Newton 常数 $G_N$，具有质量平方倒数的量纲，
\begin{equation}
  G_N \sim \frac{1}{M_{Pl}^2},
  \qquad M_{Pl} \sim 10^{19}\ \mathrm{GeV} ,
  \label{eq:22.19}
\end{equation}
其中 $M_{Pl}$ 是 Planck 质量。按照第~10 章的分析，具有负质量量纲耦合的理论是不可重整的：发散无法被有限数目的抵消项除去。因此，作为量子场论来处理的引力量子理论具有无穷多个紫外发散。

尽管存在这一困难，仍然可以在引力相互作用很弱的区域中利用量子场论作出预言。由于耦合如此之小，量子引力的首阶效应是可计算的：它们是对引力势的如下形式的修正
\begin{equation}
  V(r) = -\frac{G_N m_1 m_2}{r}\ \Bigl[ 1 + c\,
    \frac{G_N\hbar}{r^2} + \cdots \Bigr] .
  \label{eq:22.20}
\end{equation}
然而，这些修正太小，在任何可预见的实验中都无法测量。

量子引力问题的解决要求在 Planck 标度处有新的物理。一个很有前途的途径是弦理论，其中基本对象不是点粒子，而是一维延展对象。在弦理论中，引力子自然地作为弦的一种激发出现，而点粒子场论的紫外发散被弦的延展性质所软化。在低能下，弦理论重现量子场论的预言，但在量级为 Planck 质量的能量下，它的结构完全不同。

对量子引力的探索仍然是理论物理学的重大前沿之一。这是一个结合了量子力学、广义相对论和量子场论最深刻思想的问题，它的解决将需要根本性的新概念。我们希望本书中发展的量子场论工具，将为解决这些问题的下一代物理学家提供基础。

对量子场论的研究远未结束。本书发展的方法在描述基本粒子行为方面已被证明极其成功，它们仍然是探索物理学前沿不可或缺的工具。在这最后一章中，我们试图传达当前研究前沿的某种激动与挑战。那些仍然存在的问题——希格斯扇区的本质、费米子质量与混合的起源、宇宙学常数问题、引力的量子化——是全部科学中最深刻的问题之列。量子场论以其数学优雅与实验成功的非凡结合，必定在解决这些问题中发挥核心作用。我们希望本书已经给予你加入这一事业的工具和灵感。

\problems

\textbf{22.1} \emph{格点规范理论作用量。} 在本习题中，我们研究格点规范理论的 Wilson 作用量。证明在连续统极限下，格点规范理论的 Wilson 作用量约化为标准的 Yang-Mills 作用量，并计算首阶修正。

\textbf{22.2} \emph{强耦合格点规范理论中的禁闭。} 在本习题中，我们推导格点规范理论强耦合极限下 Wilson 圈的面积律。对 Wilson 圈期望值作强耦合展开，并证明它以 $e^{-\sigma A}$ 衰减，其中弦张力是可计算的。

\textbf{22.3} \emph{大统一理论中的重子数破坏。} 在本习题中，我们计算最小 $SU(5)$ 大统一理论中质子衰变的速率。证明衰变 $p\to e^{+}\pi^0$ 由重 $X$ 和 $Y$ 规范玻色子的交换所介导，并以大统一标度计算质子的寿命。
